\documentclass[aspectratio=169,usepdftitle=true,11pt,ngerman,t]{beamer}

\newcommand\lqfirstname{Name}
\newcommand\lqfullname{Name Last}
\newcommand\lqemail{name.last}
\errorcontextlines999999

% base
\usepackage[utf8]{inputenc}
\usepackage[T1]{fontenc}
\usepackage{lmodern}
\usepackage{microtype}
\usepackage[main=ngerman]{babel}
\usepackage{csquotes}

% graphics
\usepackage[prefix=]{xcolor-material}
\usepackage{graphics}
\usepackage{tcolorbox}
\usepackage{twemojis}
\usepackage{tikz}
\usetikzlibrary{
    positioning,
    calc,
    tikzmark,
    circuits.logic.IEC
}

% math
\usepackage{amsmath,amssymb}
\usepackage{mathtools}
\usepackage{array}
\usepackage{booktabs}
\usepackage{cancel}
\usepackage{ulem}

% flo
\usepackage{fancyqr}
\usepackage[addons]{color-palettes}
\usepackage[encoding,noloadlangs,cpalette,numinpar,fakeminted]{sopra-listings}
\usepackage{code-animation}
\usetheme[theorems]{colorful-dream}
\theoremstyle{exercise}
\newtheorem*{exercise*}{Aufgabe}
\theoremstyle{solve}
\newtheorem*{solve*}{Lösung}

% tools
\usepackage{algpseudocode}
\usepackage{dirtree}

% warnings
\usepackage{silence}
\WarningFilter{latexfont}{Some font shapes were not available}
\WarningFilter{latexfont}{Font shape}
\WarningFilter{pgf}{Returning node center instead of a point on node border}
\let\origenddocument\enddocument\renewcommand{\enddocument}{\hfuzz=10.95pt\origenddocument}
% title image: box
% ([vshift], inverse scale, width, content)
\newcommand\titleimagebox[4][0pt]{%
\newlength\tibheight\setlength\tibheight{4cm}
\newlength\tibtY\setlength\tibtY{-6em}
\titleimage{\begin{minipage}[c][#2\tibheight]{#3}\vspace*{#2\tibtY}%
    \vspace*{#1}\begin{tcolorbox}[boxrule=0.1pt]%
    \centering#4% TODO: bug: not centered if too wide
    \end{tcolorbox}%
\end{minipage}}}

% outro image: text
% (text)
\newcommand\outroimagetext[1]{\outro{Karlsruhe, \today%
\begin{tikzpicture}[remember picture,overlay]
    \node[align=center] at (current page.center) {#1};
\end{tikzpicture}}}% chktex 31

% reset counters
\newcommand\resetframecounters{\addtocounter{exercise}{-1}\addtocounter{solve}{-1}}
\newcommand\resetex{\addtocounter{exercise}{-1}}
\newcommand\resetsolve{\addtocounter{solve}{-1}}

% small footnote
\newcommand{\nicefootnote}[1]{\tikz[overlay,remember picture]{\node[anchor=south,text width=\dimexpr\textwidth+2em\relax,align=center,above=1em of current page.south] (nicefootnote) {\hrule\vspace{1pt}\footnotesize#1};}}

% kv
% ([draw], variables, values)
% {currently only works correctly for 4 variables}
\newcommand{\lqkv}[3][]{
    \begin{tikzpicture}
        \foreach \v [count=\k from 1] in {#2} {\global\let\Vdim=\k}
        \pgfmathsetmacro\rows{2^ceil(\Vdim/2)}
        \pgfmathsetmacro\cols{2^int(\Vdim/2)}

        \foreach \v [count=\k from 1] in {#2} {
            \pgfmathsetmacro\shift{0.5+floor(\k*ceil(\Vdim/2)/2)}

            \ifcase\k\relax
            \or
                \draw (\shift-1,0.75) -- (\shift+1,0.75);
                \node[fill=white] at (\shift,0.75) {\v};
            \or
                \draw (-0.75,-\shift) -- (-0.75,-\shift+2);
                \node[fill=white] at (-0.75,-\shift+1) {\v};
            \or
                \draw (\shift-2,-\cols+0.25) -- (\shift,-\cols+0.25);
                \node[fill=white] at (\shift-1,-\cols+0.25) {\v};
            \or
                \draw (\cols-0.25,-\shift+1) -- (\cols-0.25,-\shift+3);
                \node[fill=white] at (\cols-0.25,-\shift+2) {\v};
            \fi
        }
      
        \foreach \v [count=\k from 0] in {#3} {
            \pgfmathsetmacro\x{mod(\k,\rows)}
            \pgfmathsetmacro\y{int(\k/\rows)}
            \node[draw, minimum width=1cm, minimum height=1cm, align=center] at (\x,-\y) {\v};
        }

        \begin{scope}
            \clip (-0.5,0.5) rectangle (3.5,-3.5);
            \foreach \r in {#1} {\draw[rounded corners=1ex] \r;}
        \end{scope}

        \global\let\Vdim\undefined
    \end{tikzpicture}
}
\title{Digitaltechnik}
\institute{Karlsruher Institut für Technologie}

\author{\lqfullname}
\email{\lqemail @student.kit.edu}% chktex 1
% credit flo
\addtobeamertemplate{title page}{}{%
\scriptsize\color{gray}\begin{tikzpicture}[overlay, remember picture]
    \node[above left = 1em and 1em of current page.south east,align=center] (n0) {\LaTeX-Vorlage von\\\href{https://github.com/EagleoutIce}{Florian Sihler}};
    \node[above = 0cm of n0] (n1) {\fancyqr[size=1.5cm,tight]{https://eagleoutice.github.io/beamer-themes/\#colorful-dream}};
\end{tikzpicture}}% chktex 31
\usepackage{tikz}
\usetikzlibrary{fadings}

% cursed braces
\makeatletter
\newcommand{\cb@dx}[1]{%
    \ifnum`#1=`(  .0652em\else% chktex 37
    \ifnum`#1=`) -.0652em\else
    \ifnum`#1=`<  .0652em\else
    \ifnum`#1=`> -.0652em\else
    \ifnum`#1=`B  .0113em\else
    0em\fi\fi\fi\fi\fi
}
\newcommand{\cb@r}[1]{%
    \ifnum`#1=`< \langle\else
    \ifnum`#1=`> \rangle\else
    \ifnum`#1=`& \backslash\else
    \ifnum`#1=`= \forall\else% im not sure about this
    #1\fi\fi\fi\fi
}
\newcommand{\cb@clip}[4][]{\scoped{%
    \clip(-.3117em,#2) rectangle (.3117em,#3);%
    \node[#1] at (\cb@dx{#4},0em) {\(\cb@r{#4}\)};% chktex 1
}}
\newcommand{\cb@ht}[1][]{\cb@clip[#1]{0em}{.5669em}}
\newcommand{\cb@hb}[1][]{\cb@clip[#1]{-.5669em}{0em}}
\newcommand{\cb@tt}[1][]{\cb@clip[#1]{.1418em}{.5669em}}
\newcommand{\cb@tm}[1][]{\cb@clip[#1]{-.1418em}{.1418em}}
\newcommand{\cb@tb}[1][]{\cb@clip[#1]{-.5669em}{-.1418em}}
\newcommand{\cb@split}[4][]{%
    \ifx\relax#4\relax
        \cb@ht[#1]{#2}\cb@hb[#1]{#3}%
    \else
        \cb@tt[#1]{#2}\cb@tm[#1]{#3}\cb@tb[#1]{#4}%
    \fi
}
\newcommand{\cb}[2][]{{\tikz[baseline=-.8ex]{\cb@split[#1]#2\relax}}}
\newcommand{\cblogo}{\cb[inner sep=0pt,shading=axis,left color=orange,right color=violet,shading angle=0]{cB}}
\makeatother% chktex-file 9 chktex-file 15 chktex-file 17 chktex-file 21 chktex-file 36

\usepackage{textcomp}

%%%%%%%%%%%%%%%%%%%%%%%%%%%%%%%%%%%%%%%%%%%%%%%%%%%%%%%%%%%%%%%%

\subtitle{Tutorium 4}
\date{09. Juni 2026}

\titleimagebox{0.5}{4cm}{\cb{\{|\{}Made with \cblogo\cb{\}|\}}}

\outroimagetext{\begin{tabular}{c}
    \cb{II}\cb{cc}\cb{hh}\cb{~~}\cb{bl}\cb{ii}\cb{ne}\cb{~b}\cb{ue}\cb{n~}\cb{tm}\cb{ee}\cb{ri}\cb{bn}\cb{ee}\cb{zn}\cb{a~}\cb{hJ}\cb{lo}\cb{tb}\cb{!!}\cb{*~}\\
    \TINY\rlap{\qquad\qquad\qquad\qquad\qquad\qquad\qquad\qquad\qquad\qquad\qquad\qquad\cb{<\}<}* für die Zeit die ich hiermit verschwendet habe\cb{>\{>}}
\end{tabular}}

%%%%%%%%%%%%%%%%%%%%%%%%%%%%%%%%%%%%%%%%%%%%%%%%%%%%%%%%%%%%%%%%

\begin{document}

\section[\cb{&ao}, \cb{g==} und \cb{yzx}]{Klammern}

\subsection{\texorpdfstring{\textbf{c}ursed \textbf{B}rackets}{cursed Brackets}}
\begin{frame}{Statt Sommerurlaub}
\centering
\hspace*{-.5cm}\begin{minipage}{2\paperwidth}
\color{gray}
\cb{(((}\cb{(()}\cb{(([}\cb{((]}\cb{((\{}\cb{((\}}\cb{((<}\cb{((>}\cb{((|}\cb{((&}\cb{((/}\hspace{1em}\cb{()(}\cb{())}\cb{()[}\cb{()]}\cb{()\{}\cb{()\}}\cb{()<}\cb{()>}\cb{()|}\cb{()&}\cb{()/}\hspace{1em}\cb{([(}\cb{([)}\cb{([[}\cb{([]}\cb{([\{}\cb{([\}}\cb{([<}\cb{([>}\cb{([|}\cb{([&}\cb{([/}\hspace{1em}\cb{(](}\cb{(])}\cb{(][}\cb{(]]}\cb{(]\{}\cb{(]\}}\cb{(]<}\cb{(]>}\cb{(]|}\cb{(]&}\cb{(]/}\hspace{1em}\cb{(\{(}\cb{(\{)}\cb{(\{[}\cb{(\{]}\cb{(\{\{}\cb{(\{\}}\cb{(\{<}\cb{(\{>}\cb{(\{|}\cb{(\{&}\cb{(\{/}\\\cb{(\}(}\cb{(\})}\cb{(\}[}\cb{(\}]}\cb{(\}\{}\cb{(\}\}}\cb{(\}<}\cb{(\}>}\cb{(\}|}\cb{(\}&}\cb{(\}/}\hspace{1em}\cb{(<(}\cb{(<)}\cb{(<[}\cb{(<]}\cb{(<\{}\cb{(<\}}\cb{(<<}\cb{(<>}\cb{(<|}\cb{(<&}\cb{(</}\hspace{1em}\cb{(>(}\cb{(>)}\cb{(>[}\cb{(>]}\cb{(>\{}\cb{(>\}}\cb{(><}\cb{(>>}\cb{(>|}\cb{(>&}\cb{(>/}\hspace{1em}\cb{(|(}\cb{(|)}\cb{(|[}\cb{(|]}\cb{(|\{}\cb{(|\}}\cb{(|<}\cb{(|>}\cb{(||}\cb{(|&}\cb{(|/}\hspace{1em}\cb{(&(}\cb{(&)}\cb{(&[}\cb{(&]}\cb{(&\{}\cb{(&\}}\cb{(&<}\cb{(&>}\cb{(&|}\cb{(&&}\cb{(&/}\\\cb{(/(}\cb{(/)}\cb{(/[}\cb{(/]}\cb{(/\{}\cb{(/\}}\cb{(/<}\cb{(/>}\cb{(/|}\cb{(/&}\cb{(//}\hspace{1em}\cb{)((}\cb{)()}\cb{)([}\cb{)(]}\cb{)(\{}\cb{)(\}}\cb{)(<}\cb{)(>}\cb{)(|}\cb{)(&}\cb{)(/}\hspace{1em}\cb{))(}\cb{)))}\cb{))[}\cb{))]}\cb{))\{}\cb{))\}}\cb{))<}\cb{))>}\cb{))|}\cb{))&}\cb{))/}\hspace{1em}\cb{)[(}\cb{)[)}\cb{)[[}\cb{)[]}\cb{)[\{}\cb{)[\}}\cb{)[<}\cb{)[>}\cb{)[|}\cb{)[&}\cb{)[/}\hspace{1em}\cb{)](}\cb{)])}\cb{)][}\cb{)]]}\cb{)]\{}\cb{)]\}}\cb{)]<}\cb{)]>}\cb{)]|}\cb{)]&}\cb{)]/}\\\cb{)\{(}\cb{)\{)}\cb{)\{[}\cb{)\{]}\cb{)\{\{}\cb{)\{\}}\cb{)\{<}\cb{)\{>}\cb{)\{|}\cb{)\{&}\cb{)\{/}\hspace{1em}\cb{)\}(}\cb{)\})}\cb{)\}[}\cb{)\}]}\cb{)\}\{}\cb{)\}\}}\cb{)\}<}\cb{)\}>}\cb{)\}|}\cb{)\}&}\cb{)\}/}\hspace{1em}\cb{)<(}\cb{)<)}\cb{)<[}\cb{)<]}\cb{)<\{}\cb{)<\}}\cb{)<<}\cb{)<>}\cb{)<|}\cb{)<&}\cb{)</}\hspace{1em}\cb{)>(}\cb{)>)}\cb{)>[}\cb{)>]}\cb{)>\{}\cb{)>\}}\cb{)><}\cb{)>>}\cb{)>|}\cb{)>&}\cb{)>/}\hspace{1em}\cb{)|(}\cb{)|)}\cb{)|[}\cb{)|]}\cb{)|\{}\cb{)|\}}\cb{)|<}\cb{)|>}\cb{)||}\cb{)|&}\cb{)|/}\\\cb{)&(}\cb{)&)}\cb{)&[}\cb{)&]}\cb{)&\{}\cb{)&\}}\cb{)&<}\cb{)&>}\cb{)&|}\cb{)&&}\cb{)&/}\hspace{1em}\cb{)/(}\cb{)/)}\cb{)/[}\cb{)/]}\cb{)/\{}\cb{)/\}}\cb{)/<}\cb{)/>}\cb{)/|}\cb{)/&}\cb{)//}\hspace{1em}\cb{[((}\cb{[()}\cb{[([}\cb{[(]}\cb{[(\{}\cb{[(\}}\cb{[(<}\cb{[(>}\cb{[(|}\cb{[(&}\cb{[(/}\hspace{1em}\cb{[)(}\cb{[))}\cb{[)[}\cb{[)]}\cb{[)\{}\cb{[)\}}\cb{[)<}\cb{[)>}\cb{[)|}\cb{[)&}\cb{[)/}\hspace{1em}\cb{[[(}\cb{[[)}\cb{[[[}\cb{[[]}\cb{[[\{}\cb{[[\}}\cb{[[<}\cb{[[>}\cb{[[|}\cb{[[&}\cb{[[/}\\\cb{[](}\cb{[])}\cb{[][}\cb{[]]}\cb{[]\{}\cb{[]\}}\cb{[]<}\cb{[]>}\cb{[]|}\cb{[]&}\cb{[]/}\hspace{1em}\cb{[\{(}\cb{[\{)}\cb{[\{[}\cb{[\{]}\cb{[\{\{}\cb{[\{\}}\cb{[\{<}\cb{[\{>}\cb{[\{|}\cb{[\{&}\cb{[\{/}\hspace{1em}\cb{[\}(}\cb{[\})}\cb{[\}[}\cb{[\}]}\cb{[\}\{}\cb{[\}\}}\cb{[\}<}\cb{[\}>}\cb{[\}|}\cb{[\}&}\cb{[\}/}\hspace{1em}\cb{[<(}\cb{[<)}\cb{[<[}\cb{[<]}\cb{[<\{}\cb{[<\}}\cb{[<<}\cb{[<>}\cb{[<|}\cb{[<&}\cb{[</}\hspace{1em}\cb{[>(}\cb{[>)}\cb{[>[}\cb{[>]}\cb{[>\{}\cb{[>\}}\cb{[><}\cb{[>>}\cb{[>|}\cb{[>&}\cb{[>/}\\\cb{[|(}\cb{[|)}\cb{[|[}\cb{[|]}\cb{[|\{}\cb{[|\}}\cb{[|<}\cb{[|>}\cb{[||}\cb{[|&}\cb{[|/}\hspace{1em}\cb{[&(}\cb{[&)}\cb{[&[}\cb{[&]}\cb{[&\{}\cb{[&\}}\cb{[&<}\cb{[&>}\cb{[&|}\cb{[&&}\cb{[&/}\hspace{1em}\cb{[/(}\cb{[/)}\cb{[/[}\cb{[/]}\cb{[/\{}\cb{[/\}}\cb{[/<}\cb{[/>}\cb{[/|}\cb{[/&}\cb{[//}\hspace{1em}\cb{]((}\cb{]()}\cb{]([}\cb{](]}\cb{](\{}\cb{](\}}\cb{](<}\cb{](>}\cb{](|}\cb{](&}\cb{](/}\hspace{1em}\cb{])(}\cb{]))}\cb{])[}\cb{])]}\cb{])\{}\cb{])\}}\cb{])<}\cb{])>}\cb{])|}\cb{])&}\cb{])/}\\\cb{][(}\cb{][)}\cb{][[}\cb{][]}\cb{][\{}\cb{][\}}\cb{][<}\cb{][>}\cb{][|}\cb{][&}\cb{][/}\hspace{1em}\cb{]](}\cb{]])}\cb{]][}\cb{]]]}\cb{]]\{}\cb{]]\}}\cb{]]<}\cb{]]>}\cb{]]|}\cb{]]&}\cb{]]/}\hspace{1em}\cb{]\{(}\cb{]\{)}\cb{]\{[}\cb{]\{]}\cb{]\{\{}\cb{]\{\}}\cb{]\{<}\cb{]\{>}\cb{]\{|}\cb{]\{&}\cb{]\{/}\hspace{1em}\cb{]\}(}\cb{]\})}\cb{]\}[}\cb{]\}]}\cb{]\}\{}\cb{]\}\}}\cb{]\}<}\cb{]\}>}\cb{]\}|}\cb{]\}&}\cb{]\}/}\hspace{1em}\cb{]<(}\cb{]<)}\cb{]<[}\cb{]<]}\cb{]<\{}\cb{]<\}}\cb{]<<}\cb{]<>}\cb{]<|}\cb{]<&}\cb{]</}% chktex 9
\end{minipage}\par
\vspace*{1cm}
\Large Was habe ich erschaffen\cb{S7?}\cb{o2!}\cb{9Z:}
\end{frame}

\section[Nächste Woche dann sogar mal korrigiert!]{Übungsblatt}

\subsection{Aufgabe 1}
\begin{frame}{Shocking!}\onslide<+->
    \begin{tikzpicture}[remember picture,overlay]
        \node (n0) at ([xshift=9.5em,yshift=-1.4em]current page.north west) {\usebeamerfont{frametitle}\texttwemoji{scream}};
        \onslide<+->\node[below right=1em of n0] (n1) {\tiny Warum darf ich keinen Emoji in den Titel packen, sondern muss ihn selber zeichnen? Das ist ja voll nervig und ich schreibe diesen Text jetzt einfach nur länger, damit er definitiv aus der Folie herausläuft, weil ich das witzig finde. Haha.};
        \draw[->,bend left] (n1.west) to (n0.south);
        \draw[->,bend left,white,shorten <=15pt] (n1.west) to (n0.south);
    \end{tikzpicture}
    \only<+->{
    \begin{exercise}
        Gegeben sei die Boolesche Funktion \(f(c,b,a)=(b\wedge a)\leftrightarrow\overline{((c\vee a)\wedge (b\nleftrightarrow a))}\).
    \end{exercise}
    \onslide<+->
    \begin{solve}[\only<.-.(8)>{Shannonscher Entwicklungssatz}\only<.(9)->{Funktionstabelle} \only<.(10)>{und DNF}\only<.(11)>{und KNF}]
        \begin{itemize}\color<.(9)->{lightgray}
            \item<+-> Entwicklung nach \(c\): \only<+>{\((c\wedge((b\wedge a)\leftrightarrow\overline{((1\vee a)\wedge (b\nleftrightarrow a))}))\vee\null\)\\\hspace{10.2em}\((\overline{c}\wedge((b\wedge a)\leftrightarrow\overline{((0\vee a)\wedge (b\nleftrightarrow a))}))\)}
            \only<+->{\((c\wedge((b\wedge a)\leftrightarrow\overline{(b\nleftrightarrow a)}))\vee\null\)\\\hspace{10.5em}\((\overline{c}\wedge((b\wedge a)\leftrightarrow\overline{(a\wedge (b\nleftrightarrow a))}))\)}
            \item<+-> Entwicklung nach \(b\): \onslide<+->\((b\wedge((c\wedge(a\leftrightarrow a))\vee(\overline{c}\wedge(a\leftrightarrow\overline{(a\wedge \overline{a})}))))\vee(\overline{b}\wedge((c\wedge a)\vee(\overline{c}\wedge a\wedge a)))\)
            \item<+-> Entwicklung nach \(a\): \onslide<+->\((a\wedge((b\wedge(c\vee\overline{c}))\vee(\overline{b}\wedge(c\vee\overline{c}))))\vee(\overline{a}\wedge b\wedge c)\)
            \item<+-> \usebeamercolor[fg]{normal text}\(f(c,b,a)=\overline{a}bc\vee abc\vee ab\overline{c}\vee a\overline{b}c\vee a\overline{b}\overline{c}\)
        \end{itemize}
    \end{solve}
    \onslide<+->\only<.>{\addtocounter{framenumber}{1}}
    \begin{tikzpicture}[remember picture,overlay]
        \node[fill=white,anchor=west] (tab12) at ([xshift=4em,yshift=-3.5em]current page.center) {\begin{tabular}{|ccc|c|}
            \hline
            $c$ & $b$ & $a$ & $f(c,b,a)$ \\ \hline
            0 & 0 & 0 & \textcolor<.(2)>{blue}{0} \\ \hline
            0 & 0 & 1 & \textcolor<.(1)>{red}{1} \\ \hline
            0 & 1 & 0 & \textcolor<.(2)>{blue}{0} \\ \hline
            0 & 1 & 1 & \textcolor<.(1)>{red}{1} \\ \hline
            1 & 0 & 0 & \textcolor<.(2)>{blue}{0} \\ \hline
            1 & 0 & 1 & \textcolor<.(1)>{red}{1} \\ \hline
            1 & 1 & 0 & \textcolor<.(1)>{red}{1} \\ \hline
            1 & 1 & 1 & \textcolor<.(1)>{red}{1} \\ \hline
           \end{tabular}};
        \only<+>{\node[below=0em of tab12,anchor=north] {\color{red}\(f(c,b,a)=\T{MINt(1,3,5,6,7)}\)};}
        \only<+>{\node[below=0em of tab12,anchor=north] {\color{blue}\(f(c,b,a)=\T{MAXt(0,2,4)}\)};}
    \end{tikzpicture}}
\end{frame}
\resetframecounters%

\subsection{Aufgabe 2}
\begin{frame}{Was ist Was?}\onslide<+->
    \begin{exercise}
        Bestimmen Sie für jeden Term in der Tabelle, ob es sich um einen Min- oder Maxterm oder die disjunktive oder konjunktive Normalform oder Minimalform handelt.
    \end{exercise}
    \begin{solve}
    \begin{itemize}[<+->]
        \item \(c\overline ba\vee \overline c\overline b\overline a\): \onslide<+-> DNF, DMF
        \item \(\overline c\vee b\vee a\): \onslide<+-> Maxterm, KNF, DMF, KMF
        \item \(c\overline ba\): \onslide<+-> Minterm, DNF, DMF, KMF
        \item \(c\vee \overline{(\overline b\wedge\overline a)}\): \onslide<+->-
        \item \(1\): \onslide<+-> KNF, DMF, KMF
        \item \(a\vee b\): \onslide<+-> DMF, KMF
    \end{itemize}
    \end{solve}
\end{frame}

\subsection{Aufgabe 3}
\begin{frame}[noframenumbering]{ChatGPT, bitte zeichne mir dieses Schaltnetz in tikz!}
    \begin{tikzpicture}[circuit logic IEC,
        every circuit symbol/.style={scale=1}]
        % Inputs
        \node (a) at (0,2) {$a$};
        \node (b) at (0,0) {$b$};
        
        % ======================
        % Middle chain: a,b -> =1 -> 1 -> OR (top)
        % ======================
        \node[xor gate IEC, draw, right=2cm of a] (xorTop) {};
        \node[buffer gate IEC, draw, right=2cm of xorTop] (bufTop) {};
        \node[or gate IEC, draw, right=2cm of bufTop] (orTop) {};
        
        % wires to xorTop
        \draw (a.east) |- (xorTop.input 1);
        \draw (b.east) |- (xorTop.input 2);
        
        % xorTop -> buffer -> orTop
        \draw (xorTop.output) -- (bufTop.input);
        \draw (bufTop.output) -- (orTop.input 2);
        
        % direct a to orTop (top wire)
        \draw (a.east) -- ++(6,0) |- (orTop.input 1);
        
        % ======================
        % Bottom chain: a,b -> =1 -> OR (bottom)
        % ======================
        \node[xor gate IEC, draw, right=2cm of b, yshift=-3cm] (xorBot) {};
        \node[or gate IEC, draw, right=4cm of xorBot] (orBot) {};
        
        % inputs to bottom XOR
        \draw (a.east) |- (xorBot.input 1);
        \draw (b.east) |- (xorBot.input 2);
        
        % XOR -> OR bottom
        \draw (xorBot.output) -- (orBot.input 1);
        
        % direct a to bottom OR
        \draw (a.east) -- ++(6,-3) |- (orBot.input 2);
        
        % ======================
        % Final OR
        % ======================
        \node[or gate IEC, draw, right=3cm of orTop, yshift=-1.5cm] (orFinal) {};
        
        \draw (orTop.output) -- (orFinal.input 1);
        \draw (orBot.output) -- (orFinal.input 2);
        
        % Output
        \draw (orFinal.output) -- ++(1,0) node[right] {$y$};
    \end{tikzpicture}
\end{frame}
\resetframecounters%
\begin{frame}{Dann halt so\dots}\onslide<+->
    \begin{enumerate}[<+->]
        \item \(y=(\overline{(a\nleftrightarrow b)}\overline\vee a)\vee((a\nleftrightarrow b)\vee b)\)
        \item \(y=a\overline b\vee b\)
        \item Tafel
    \end{enumerate}
\end{frame}

\begin{frame}{es ist 15:33}\onslide<+->
    Daher an der Tafel

    \onslide<+->\scriptsize(Jetzt ist es 15:38)
\end{frame}

\end{document}